% Source: Independent modeling summary; current handbook conventions.
\subsection{状态、初始化与转移顺序}
\label{dp:basics}
\subsubsection*{原理}

写一个 DP 时，先依次回答：状态 $f$ 的下标各代表什么；哪些状态可达及初值是什么；最后一步可以怎样选择、转移来自哪里；按什么顺序计算；最后读哪个状态或怎样汇总。代码中的数组、循环与答案应能逐项对应这五点。

状态应保留未来决策所需的全部信息，并明确记录最优值、方案数还是概率。按依赖顺序计算；记忆化搜索也要有严格减小的规模，避免循环依赖。

求最大值时，不可达状态设为 $-INF$，求最小值时设为 $INF$；每次转移先判断前驱是否可达，再做加法。计数从“空方案有一种”开始，概率从“初始事件概率为 1”开始，不能沿用最值 DP 的初始化。

滚动数组只压掉不再使用的层，不改变原转移。若压成一维，必须判断每次读取的是上一层还是当前层，不能先套正序或倒序的经验。每次重新做一组数据都重建状态；保留路径时，记录的前驱也必须属于正确的层。

本章以普通数组、循环和递归函数展示这些步骤，不要求为每个 DP 建立结构体。具体题目通常要修改状态和转移；凸包等有独立维护操作的数据结构仍单独封装。

复杂度按\textbf{状态数乘每个状态的转移数}估算，并把预处理、状态压缩的 $2^n$ 和容器开销计入。优化前先保留直接转移，后续用它核对新算法。
